\section{可审查的核心规范}\label{sec:appendix:specification}

本附录给出不依赖源码文件名即可检查的规范性状态机、权限、事件、恢复、
归责、影响修复和最终论文验证。未列出的转换与权限一律拒绝。

\subsection{完整的 T1--T21 组件生命周期转换表}\label{sec:appendix:transitions}

十个状态为候选、已验证、影子、试用活动、活动、警告、观察、隔离、退役
和封禁。所有规则都在边界事务中提交；T16 会立即形成紧急边界。

\small
\begin{longtable}{@{}p{0.045\textwidth}p{0.17\textwidth}p{0.31\textwidth}p{0.40\textwidth}@{}}
  \caption{完整组件生命周期转换表。未列出的边均为非法。}
  \label{tab:component-transitions}\\
  \toprule
  \textbf{规则} & \textbf{状态变化} & \textbf{触发条件} & \textbf{附加条件与效果} \\
  \midrule
  \endfirsthead
  \toprule
  \textbf{规则} & \textbf{状态变化} & \textbf{触发条件} & \textbf{附加条件与效果} \\
  \midrule
  \endhead
  T1 & 候选$\rightarrow$已验证 & 候选验证通过 & 未超过候选上限且包含宪制约束。\\
  T2 & 候选$\rightarrow$退役 & 验证失败或过期 & 仅拒绝从未服务的候选。\\
  T3 & 已验证$\rightarrow$影子 & 重放与锚点评估通过 & 存在影子名额。\\
  T4 & 影子$\rightarrow$已验证 & 影子样本不足 & 重试次数仍在上限内。\\
  T5 & 影子$\rightarrow$退役 & 影子质量低于阈值 & 仅拒绝从未服务的候选。\\
  T6 & 影子$\rightarrow$试用活动 & 一致性达到阈值 &
    角色存在空缺；每角色每边界最多采用一个。\\
  T7 & 试用活动$\rightarrow$活动 & 试用期无违规 & 满足最短试用期。\\
  T8 & 试用活动$\rightarrow$隔离 & 弹劾成立或可靠性崩溃 & 在普通边界执行。\\
  T9 & 活动$\rightarrow$警告 & 可靠性警告或低严重度攻击 & 目标继续服务。\\
  T10 & 活动$\rightarrow$观察 & 警告重复或中等攻击 & 目标继续服务。\\
  T11 & 活动$\rightarrow$隔离 & 弹劾成立或严重攻击 & 指定后备角色。\\
  T12 & 警告$\rightarrow$活动 & 一个干净纪元 & 解除警告。\\
  T13 & 警告$\rightarrow$观察 & 记忆期内再次警告 & 目标继续服务。\\
  T14 & 观察$\rightarrow$活动 & 完成所需干净纪元 & 满足最短观察期。\\
  T15 & 观察$\rightarrow$隔离 & 持续退化或弹劾成立 & 在普通边界执行。\\
  T16 & 服务状态$\rightarrow$隔离 & 固定内核的严重违规 &
    仅允许从试用活动、活动、警告或观察出发；必须附违规证据、只制裁
    一个目标，并强制关闭/开启纪元。\\
  T17 & 隔离$\rightarrow$退役 & 退役裁决确认 & 需要重放确认并追加退役记录。\\
  T18 & 隔离$\rightarrow$试用活动 & 免责 & 从试用期重新进入。\\
  T19 & 退役$\rightarrow$封禁 & 污染或严重发现 & 封禁为吸收状态。\\
  T20 & 活动$\rightarrow$退役 & 已采用的后继替代现任 &
    仅效用策略角色可用，并追加退役记录。\\
  T21 & 活动$\rightarrow$影子 & 已采用的后继替代现任 &
    仅论文写作者/审稿者可用，旧版本作为可恢复后备保留。\\
  \bottomrule
\end{longtable}
\normalsize

T16 不再在同一纪元内改变服务集合。它强制关闭当前纪元，并在一个事务中
提交隔离、旧纪元关闭和具有新指纹的新纪元开启。原效用策略保持原样冻结，
紧急事务不采用后继。因此 T16 前后的样本具有不同制度标识。

\paragraph{纪元身份。}
复合纪元指纹包含两个分别报告的摘要。\emph{政策指纹}提交活动组件与提示词
版本、全部已注册提示词摘要、效用策略、注册表版本、宪制摘要、转换阈值、
预算及其他已解析治理设置。\emph{执行指纹}提交声明的模型后端与标识、温度、
搜索与评估设置、禁用工具、技能集合，以及模型修订、工具包、环境和数据快照的
显式固定值。
若外部修订不可得，就记录为“未解析”，而不默认为已固定。因此政策指纹相等
只表示声明的治理政策相等；完整执行指纹相等表示声明的执行依赖相等。二者
都不能证明提供者内部不透明权重相等。

\paragraph{边界解析顺序。}
解析器先冻结报告、候选评估、注册表和设置，再按标识顺序处理组件，然后
按标识顺序处理没有配对组件的提示词候选。组件处分建立角色空缺；每个对象
最多走一条生命周期边。T6 还要求空缺，并限制每角色每边界最多一次采用。
合格的效用策略后继把现任 T20 退役与后继 T6 采用放入同一事务；论文角色
后继则在同一事务中以 T21 把现任移回影子。事件按采用、处分、退役、封禁
的固定组顺序序列化，但事务提交前不存在可见中间状态。随后导出事后活动
映射；每个被移除的服务角色使用最近的合格旧持有者，否则使用固定基线。
内核验证后才原子提交单一边界事务。

\begin{table}[htbp]
  \centering
  \small
  \begin{tabularx}{\textwidth}{@{}p{0.34\textwidth}X@{}}
    \toprule
    \textbf{全局状态不变量} & \textbf{当前强制方式} \\
    \midrule
    每角色至多一个主要服务者 &
      确定性事后投影与 T6 角色上限 \\
    每个复合指纹只对应一个服务集合 &
      仅边界更新；T16 创建新纪元 \\
    候选权限不得超过现任 &
      采用前固定的权限子集检查 \\
    未裁决攻击不能处分或计分 &
      内核与评分路径要求验证攻击 \\
    依赖退役角色的记录不能重新入选 &
      修复闭包与资格检查 \\
    未提交事务不重放 &
      准备、成员、提交投影规则 \\
    T16 后不再生成旧纪元身份样本 &
      紧急关闭、隔离、开启事务 \\
    \bottomrule
  \end{tabularx}
  \caption{实现强制的组合状态不变量。它们有回归测试，但尚未经过模型检查
  或形式证明。}
  \label{tab:whole-state-invariants}
\end{table}

\subsection{权限矩阵}\label{sec:appendix:authority}

封闭操作集合为读取、追加、写入、调用和激活；资源为记录、注册表、审计
日志、前沿、活动提示词、退役提示词、检查点成果和候选。未列出的组合均
拒绝，任何角色都不得读取退役提示词。

\begin{table}[htbp]
  \centering
  \scriptsize
  \setlength{\tabcolsep}{3.5pt}
  \begin{tabularx}{\textwidth}{@{}p{0.19\textwidth}*{8}{>{\centering\arraybackslash}X}@{}}
    \toprule
    \textbf{角色--层级} & \textbf{记录} & \textbf{注册表} &
    \textbf{审计} & \textbf{前沿} & \textbf{活动词} &
    \textbf{退役词} & \textbf{成果} & \textbf{候选} \\
    \midrule
    可演化--制度层 & 读追调 & -- & -- & -- & 读 & -- & 读 & -- \\
    可演化--元层 & 读追 & -- & -- & -- & 读 & -- & -- & -- \\
    效用策略--制度层 & 读追 & -- & -- & -- & -- & -- & -- & -- \\
    治理角色--制度层 & 读追调 & -- & 读追 & -- & 读 & -- & 读 & -- \\
    转换器--固定层 & 读 & 读写 & 读追 & 读 & -- & -- & 读 & 激 \\
    内核/验证器--固定层 & 读 & 读 & 读 & 读 & -- & -- & 读 & -- \\
    审计日志--固定层 & -- & -- & 追 & -- & -- & -- & -- & -- \\
    \bottomrule
  \end{tabularx}
  \caption{按资源展开的完整稀疏权限矩阵。读、追、写、调、激分别表示
    读取、追加、写入、调用和激活。空白项与所有退役提示词读取均拒绝。
    实现还输出完全展开的角色--层级--操作--资源机器可读清单。}
  \label{tab:authority-matrix}
\end{table}

候选权限必须是所替代现任权限的子集。运行时的类型化构造器忽略模型文本
中的作者或角色字段，按当前组件标识和固定角色加盖。因此模型输出不能仅靠
自报身份通过应用边界。但这不是密码学或操作系统隔离：角色共享进程，没有
按角色分配的密钥或 OS 用户，也没有 IPC 沙箱。运行时、检查点权限和工具
网关属于当前可信计算基。

\begin{table}[htbp]
  \centering
  \small
  \begin{tabularx}{\textwidth}{@{}p{0.28\textwidth}XX@{}}
    \toprule
    \textbf{攻击者能力} & \textbf{当前强制} & \textbf{不保证} \\
    \midrule
    受治理模型的恶意文本 &
      类型化作者加盖、API 许可表、合法转换 &
      模型证据的语义真实性 \\
    模型角色串谋或相关故障 &
      独立记录与拒绝自我裁决 &
      统计独立性或抗串谋能力 \\
    进程内角色代码故障或被攻陷 &
      经正常网关时执行检查 &
      阻止直接内存或文件访问 \\
    检查点写入者或主机管理员 &
      无对抗性保证 &
      注册表、日志、本地固定值或可执行文件完整性 \\
    稳定模型名背后的提供者变更 &
      在执行身份中记录不完整性 &
      不透明模型权重相等 \\
    外部不可逆副作用 &
      本成果物没有固定副作用闸门 &
      设备、API 或生产数据的故障关闭安全 \\
    \bottomrule
  \end{tabularx}
  \caption{威胁能力矩阵。当前主张是在可信主机上的应用级引用监控与工作流
  完整性，不是操作系统或密码学安全边界。}
  \label{tab:threat-capability}
\end{table}

\subsection{事件封装与恢复}\label{sec:appendix:events}

新状态事件使用模式版本二，包含单调标识、类型、事务标识、载荷、前驱摘要
和时间戳。完整 SHA-256 提交规范化元组
\[
  (\textit{模式版本},\textit{事件标识},\textit{事件类型},
  \textit{事务标识},\textit{载荷},\textit{前驱摘要})
\]
因此类型、标识、事务边界、前驱、载荷或顺序变化都会改变摘要。时间戳是
摘要外的元数据。旧版一的仅载荷事件仍可读取，但新写入只用版本二。摘要
并非作者签名。

\begin{algorithm}[htbp]
\caption{确定性的恢复投影}
\begin{algorithmic}[1]
\Require 按存储顺序排列的事件 $E$，以及仍存在时的本地边界固定值 $P$
\Ensure 重建的纪元与注册表，或完整性违规
\State 检查标识严格递增
\State 重算每个版本相应的事件摘要与前驱链接
\If{$P$ 存在}
  \State 要求 $|E|\geq P.length$ 且 $P.length$ 处哈希等于 $P.head$
\EndIf
\State 要求准备、成员与提交的事务标识相同
\State 丢弃没有对应提交的末尾准备事务
\State 从空状态依次折叠已提交事件，重建纪元与注册表
\State \Return 重建状态
\end{algorithmic}
\end{algorithm}

版本二能发现重放相关字段改动、事件重排、事务边界替换、保留后继时的
中间删除，以及比仍存本地固定值更短的日志。写入方拒绝扩展损坏链；
重放仅把未提交的最后事务作为崩溃恢复忽略。它仍不能发现日志与本地固定值
一起回滚到同一旧前缀。当前实现没有
外部签名头、硬件单调计数器或独立见证者，因此整检查点回滚不在威胁模型内。

\subsection{归责、影响修复与论文验证}\label{sec:appendix:algorithms}

\paragraph{归责。}
原始攻击只针对成果。裁决后，固定的攻击类型表选择相关角色，并解析为
纪元冻结的现任。每条验证攻击只命名一个责任组件；多角色事件为每个角色
各建一条、共同引用同一裁决。七类一般研究攻击只有在目标节点带有写入一次
的生成组件、提示词和纪元来源，且与纪元冻结生成者一致时才绑定生成者。
旧记录、缺失或不匹配的来源保持无目标，不作猜测。论文过度接受绑定审稿者；
草稿同时未通过固定忠实性门时另建写作者记录。多个原因用共享同一裁决的
多条单目标记录表示。固定映射的因果准确性尚未评估；单目标形式只是审计
约束，不是因果保证，工具、数据或共享提示词原因尚不能表示。

\paragraph{影响修复。}
直接影响是带有退役提示词哈希的记录；间接影响通过反向来源引用边做广度
优先遍历。只有提案引用提案作为灵感被声明为非实质；未知边默认实质。
遍历在默认实质性深度八以后仍计算完整可达闭包；超过该深度后不再豁免
非实质边，故更深的可达记录保守失效。若可达记录数为 $V$、反向引用数为
$E$，时间与内存均为 $O(V+E)$，访问集合会终止环。生成/路由/策略的
直接影响使节点失效；审稿/对抗/辩护/裁决影响从剩余输入重算，输入不足则
失效。语义精度和过度移除率尚未评估。

\paragraph{最终论文验证。}
验证器不是通用自然语言推理器。论文前向声明主张和数值标识、公式、操作数
的节点/测量位置、单位及容差。验证器检查已执行节点与成果存在，使用封闭
公式集合重算，并按绝对或相对容差比较。摘要、结果和结论接受严格扫描；
引言、讨论和局限产生警告；相关工作、参考文献、附录和公式排除。它报告
可复现数值主张率、数值覆盖率、标识冲突、未知公式、缺少操作数、单位和
环境不匹配。本报告并非由 ARI 执行检查点生成，因此不存在有效的执行证据
主张清单，也不声称完成自我应用。用手工检查点代替不能检验所提发布路径。
非结构化非数值主张、
执行记录本身的真实性和语义改写均不在范围内。“固定”仅表示对结构化输入
不调用模型且结果确定。

\subsection{确定性与测试证据的边界}\label{sec:appendix:determinism}

本文只主张\emph{内核确定性}：相同的已保存建议报告、候选评估、注册表和
配置产生相同转换与修复；以及\emph{日志重放确定性}：相同已提交事件列
产生相同投影。不主张模型能从同一提示词重新生成相同辩护、裁决建议或审稿。

本次修订的完整套件收集 4,805 项：4,787 项通过，16 项跳过，2 项为预期
失败。直接覆盖 RQGM 和最终主张门的 1,124 项全部通过，并新增事件语义
字段、前驱、排序、事务不匹配、紧急边界指纹及权限清单完整性的负例。
无重叠分组为：主张门 64；内核等 285；对抗/治理 126；修复/净室 66；
效用演化 85；评估 157；提示词演化 93；论文角色 94；其他 154。测试数量不是性能
证据；代码覆盖率、变异测试得分、多模型检测率、成本和研究质量比较均尚未
测量。
